%============================================================
% TÓM TẮT NỘI DUNG ĐỒ ÁN
%============================================================
\chapter*{TÓM TẮT NỘI DUNG ĐỒ ÁN}
\addcontentsline{toc}{chapter}{Tóm tắt nội dung đồ án}

Ô nhiễm không khí đang là một trong những vấn đề môi trường nghiêm trọng
nhất hiện nay, đặc biệt tại các đô thị lớn như Hà Nội. Tuy nhiên, các hệ thống
quan trắc chất lượng không khí hiện có thường có chi phí triển khai cao, mật độ
trạm thưa và thiếu khả năng dự báo sớm. Đồ án trình bày việc thiết kế và triển
khai một hệ thống AIoT giám sát và dự báo chất lượng không khí theo thời gian
thực, sử dụng phần cứng ESP32 chi phí thấp kết hợp với mô hình học máy trên dữ
liệu chuỗi thời gian.

Node cảm biến sử dụng ESP32, DHT11 để đo nhiệt độ, độ ẩm và sinh dữ liệu AQI mô
phỏng nhằm hoàn thiện pipeline giám sát chất lượng không khí. Thiết bị hỗ trợ
cấu hình Wi-Fi ban đầu qua AP \texttt{Sensor-Setup-[MAC]}, đăng ký với backend
qua MQTT topic \texttt{devices/register}, nhận \texttt{device\_id} và
\texttt{api\_key}, sau đó gửi payload \texttt{temperature}, \texttt{humidity},
\texttt{aqi} định kỳ 5 giây qua topic \texttt{device/\{api\_key\}/data}.

Backend Node.js/Express tích hợp Aedes MQTT broker, PostgreSQL, JWT và Socket.io
để xác thực thiết bị, lưu dữ liệu, kiểm tra cảnh báo và đẩy sự kiện
\texttt{new\_reading} tới dashboard theo thời gian thực. Frontend được xây dựng
bằng React 19 và Vite, sử dụng Chart.js cho biểu đồ, Leaflet cho heatmap không
gian và các trang Dashboard, Devices, Forecast, Spatial Forecast, Rooms, Alerts.

Pipeline dự báo được triển khai bằng FastAPI và PyTorch. Mô hình LSTM sử dụng ba
feature đầu vào gồm nhiệt độ, độ ẩm và AQI, lookback 30 phút và forecast horizon
15 phút. Với bài toán không gian, service \texttt{spatial\_forecast.py} kết hợp
IDW để sinh node ảo ở biên phòng và Kriging để nội suy lưới 40$\times$50 điểm,
trả về dữ liệu JSON cho React Leaflet hiển thị heatmap.

Kết quả triển khai cho thấy hệ thống có thể thu thập, lưu trữ, hiển thị realtime,
dự báo và trực quan hóa phân bố không gian trong một kiến trúc hoàn chỉnh, phù
hợp cho nghiên cứu, giáo dục và mở rộng thành hệ thống giám sát môi trường thực
tế.


